\documentclass[11pt]{report}

\usepackage[margin=1in]{geometry} 
\usepackage{amsmath,amsthm,amssymb,amsfonts}

\usepackage[utf8]{inputenc}
\usepackage[T1]{fontenc}
\usepackage[spanish]{babel}


\usepackage{mathpazo,euler}
\linespread{1.05}
\usepackage{mathrsfs}

\usepackage{enumitem}
\usepackage{microtype}
\usepackage{tocbibind}

\usepackage{thmtools,xcolor}
\usepackage{fancyhdr}
\pagestyle{fancy}
\usepackage[misc]{ifsym}
\usepackage{tcolorbox}
\tcbuselibrary{theorems}

\usepackage{tikz}
\usepackage{tikz-cd}
\usetikzlibrary{arrows}
\usetikzlibrary{matrix}

\definecolor{color}{RGB}{0, 131, 126}
\addto\captionsspanish{\renewcommand{\chaptername}{Parte}}
\renewcommand\qedsymbol{$\tpaint{\small{\FilledSectioningDiamond}}$}
\declaretheoremstyle[
  headfont=\color{color}\normalfont\bfseries,
  notefont=\color{color}\normalfont\bfseries
]{colored}
\theoremstyle{colored}
\newtheorem{definition}{Definición}[section]
\newtheorem{theorem}{Teorema}[section]
\newtheorem*{theorem*}{Teorema}
\newtheorem{proposition}{Proposición}[section]
\newtheorem{corollary}{Corolario}[section]
\newtheorem{lemma}{Lema}[section]
\newtheorem{remark}{Observación}[section]
\newtheorem{example}{Ejemplo}[section]

\newcommand{\N}{\mathbb{N}}
\newcommand{\Z}{\mathbb{Z}}
\newcommand{\Q}{\mathbb{Q}}
\newcommand{\R}{\mathbb{R}}
\newcommand{\C}{\mathbb{C}}
\newcommand{\M}[2]{\mathsf{M}_{#1}#2}
\newcommand{\im}{\operatorname{im}}
\newcommand{\id}{\operatorname{id}}
\newcommand{\eps}{\varepsilon}
\newcommand{\nat}[1]{[\![#1]\!]}
\newcommand{\ord}[1]{\nat{#1}}
\newcommand{\natzero}[1]{\nat{#1}_0}
\newcommand{\ol}{\overline}
\newcommand{\cat}[1]{\mathsf{#1}}
\renewcommand{\ss}[1]{\Delta^{#1}}
\newcommand{\bss}[1]{\partial \ss{n}}
\newcommand{\horn}[2]{\Lambda^{#1}_{#2}}
\newcommand{\ordcat}{\boldsymbol{\Delta}}
\newcommand{\catlim}[2]{\underset{#1}{\operatorname{lim}}#2}
\newcommand{\catcolim}[2]{\underset{#1}{\operatorname{colim}}#2}

\newcommand{\paint}[1]{\color{color}{#1}}
\newcommand{\tpaint}[1]{\paint{\textbf{#1}}}
\newcommand{\paintline}{\begin{center}
$\paint{
\rule{400pt}{0.5pt}
}$
\vspace{10pt}
\end{center}}

%-----------------------

\title{
\LARGE{Topología Algebraica}
\\
\vspace{3pt}
\small{Primer Cuatrimestre -- 2019}
\\
\vspace{0.5pt}
\large{Examen Final}
\\
\vspace{80pt}
{\includegraphics[height=5cm]{uba2.jpg}}
\vspace{80pt}
}
\author{Guido Arnone}
\date{}
\lhead{Guido Arnone}
\rhead{Examen Final}

\begin{document}

\maketitle
\tableofcontents

\chapter{Preliminares}

\section{La categoría de ordinales finitos}

\begin{definition} Se define la \textbf{categoría $\ordcat$ de ordinales finitos} como la categoría que tiene por objetos a los conjuntos ordenados
\[
\ord{n} := \{0 < 1 < \cdots < n\}
\]
para cada $n \in \N_0$, y cuyas flechas son las funciones $f : \ord{n} \to \ord{m}$ que resultan morfismos de posets. Definimos además, para cada $\ord{n} \in \ordcat$ e $i \in \natzero{n}$, 
\begin{itemize}
\item los mapas de \textbf{cocaras},
\begin{align*}
d^i : \ord{n-&1} \to \ord{n}\\
&j \mapsto \begin{cases}
j &\text{si $j < i$}\\
j+1 &\text{si $j \geq i$}
\end{cases}
\end{align*}
y
\item los mapas de \textbf{codegeneraciones},
\begin{align*}
s^i : \ord{n+&1} \to \ord{n}\\
&j \mapsto \begin{cases}
j &\text{si $j \leq i$}\\
j-1 &\text{si $j > i$}
\end{cases}
\end{align*}
\end{itemize} 
\end{definition}

\begin{proposition} Los mapas de cocaras y codegeneraciones satisfacen las siguientes \textit{identidades cosimpliciales},
\[
\begin{cases}
d^jd^i = d^id^{j-1} &\text{si $i < j$}\\
s^jd^i = d^is^{j-1} &\text{si $i < j$}\\
s^jd^j = s^jd^{j+1} = 1 &\text{si $i < j$}\\
s^jd^i = d^{i-1}s^j &\text{si $i > j+1$}\\
s^js^i = s^is^{j+1} &\text{si $i \leq j$}\\
\end{cases}
\]
\qed \\
\end{proposition}

\begin{proposition} Toda flecha $\ord{n} \xrightarrow{f} \ord{m}$ en $\ordcat$ se puede escribir como una composición de mapas de cocaras y codegeneraciones. \\ \qed
\end{proposition}

\begin{remark} En vista de las dos proposiciones anteriores,  usando los mapas de cocaras y codegeneraciones y las identidades cosimpliciales se puede dar «una presentación de $\ordcat$ en términos de generadores y relaciones». A grandes rasgos, esto nos permitirá definir los objetos relacionados a $\ordcat$ únicamente a partir de los mapas de cocaras y codegeneraciones.
\end{remark}

\section{Conjuntos Simpliciales}

Ahora sí, pasamos a definir los conjuntos simpliciales:

\begin{definition} Un \textbf{conjunto simplcial} es un funtor $X : \ordcat^{op} \to \cat{Set}$. Concretamente, éste consiste de 
\begin{itemize}
\item[(i)] una sucesión $X_0,X_1,X_2, \dots$ de conjuntos, y \item[(ii)] para cada $n \in \N_0$ e $i \in \natzero{n}$, funciones 
$d_i : X_n \to X_{n-1}$ y $s_i : X_n \to X_{n+1}$ llamadas mapas de caras y degeneraciones respectivamente, que satisfacen las siguientes \textit{identidades simpliciales}:
\[
\begin{cases}
d_id_j = d_{j-1}d_i &\text{si $i < j$}\\
d_is_j = s_{j-1}d_i &\text{si $i < j$}\\
d_js_j = d_{j+1}s_j = 1 &\text{si $i < j$}\\
d_is_j = s_jd_{i-1} &\text{si $i > j+1$}\\
s_is_j = s_{j+1}s_i &\text{si $i \leq j$}\\
\end{cases}
\] 
\end{itemize}
\end{definition}

\begin{example}[el $n$-símplex estándar] Para cada $n \in \N_0$ tenemos un conjunto simplicial dado por el funtor $\ordcat(-,\ord{n})$. Conrectamente, para cada $m \geq 0$ definimos los conjuntos
\[
\ss{n}_m := \ordcat(\ord{m},\ord{n}) = \{f : \ord{m} \to \ord{n} : \text{ $f$ es morfismo de posets } \}
\]
y los mapas de caras y degeneraciones están dados por
\[
(d^i)^* : f \in \ordcat(\ord{m},\ord{n}) \mapsto fd^i \in \ordcat(\ord{m-1},\ord{n})
\]
y
\[
(s^i)^* : f \in \ordcat(\ord{m},\ord{n}) \mapsto fs^i \in \ordcat(\ord{m+1},\ord{n}).
\]
Llamamos a este conjunto simplicial el \textbf{n-símplex estándar} y lo notamos $\ss{n}$.
\end{example}

Una vez más interpretando a $\ord{n}$ como el $n$-símplex combinatorio ordenado, su conjunto simplicial asociado consiste de «todas las formas posibles de incluir o colapsar un $m$-símplex estándar en $\nat{n}$».

Extendiendo esta interpretación tenemos el siguiente ejemplo,

\begin{example}[complejos simpliciales ordenados] Sea $K$ un complejo simplicial equipado con una relación de orden total para sus vértices $V = \{v_i\}_{i \in I}$. Notamos a cada $n$-símplex como una $n$-upla $[v_{i_0},\dots, v_{i_n}]$ con $v_k < v_{k+1}$ para cada $k$.

Asociaremos a $K$ un conjunto simplicial, agregando como en el caso de $\ss{n}$ la noción de \textit{símlpices degenerados}. Concretamente, para cada $n \in \N_0$ definimos
\[
K_n := \left\{[v_{i_0}, \dots, v_{i_n}] : \text{ $v_{i_k} \leq v_{i_{k+1}}$ para cada $k$, y $\{v_{i_k}\}_{k = 0}^n \in K$ }\right\}.
\]

En otras palabras, el conjunto $K_n$ consiste de $n$-uplas ordenadas de vértices que forman un símplex de $K$, pero permitiendo repetición. Definimos a su vez los mapas de caras y degeneraciones como
\[
d_k[v_{i_0}, \dots, v_{i_n}] := [v_{i_0},\dots,\widehat{v_{i_k}},\dots, v_{i_n}] \quad \text{y} \quad
s_k[v_{i_0}, \dots, v_{i_n}] := [v_{i_0},\dots,v_{i_k},v_{i_k}, \dots, v_{i_n}].
\]
\end{example}

\begin{definition} Dado un $n$-śimplex estándar (como complejo simplicial), su \textbf{borde} es el complejo $\bss{n}$ dado por la unión de sus caras maximalesy su $k$-ésimo \textbf{cuerno} es el subcomplejo $\horn{n}{k}$ de $\bss{n}$ que se obtiene quitando la $k$-ésima cara maximal de $\ss{n}$, para cierto $0 \leq k \leq n$. Decimos que $\horn{n}{k}$ es un cuerno \textit{interno} si $0 < k < n$, y \textit{externo} en caso contrario.
\end{definition}

Ahora sí, veamos un primer ejemplo topológico:

\begin{example} Sea $X$ un espacio topológico. Para cada $n \in \N_0$ definimos el conjunto
\[
\mathcal{S}(X)_n := \cat{Top}(|\ss{n}|,X)
\]
de todos los $n$-simplices singulares de $X$, y las aplicaciones 
\[
d_i : \mathcal{S}(X)_n \to \mathcal{S}(X)_{n-1}, \quad
s_i : \mathcal{S}(X)_n \to \mathcal{S}(X)_{n+1}
\]
que envían un $n$-símplex singular a la restricción $d_i\sigma$ a su $i$-ésima cara y al $(n+1)$-simplex singular $s_i\sigma$ que corresponde a colapsar $|\ss{n+1}|$ a $|\ss{n}|$ a través de $|s^i|$ y luego componer con $\sigma$.

Estos conforman el conjunto simplicial $\mathcal{S}(X)$ que se conoce como el \textbf{conjunto singular} de $X$.
\end{example}

\begin{definition} Dado un conjunto simplicial $X$, definimos su \textbf{complejo de Moore} como el complejo de cadenas
\[
\cdots \to \Z X_2 \xrightarrow{\partial} \Z X_1 \xrightarrow{\partial} \Z X_0,
\]
con $\Z X_n$ el grupo abeliano libre generado por el conjunto $X_n$ y 
\[
\partial = \sum_{i=1}^n(-1)^i d_i
\]
para cada $n \geq 0$.
\end{definition}

\begin{remark} La homología singular de un espacio topológico $X$ coincide con la homología del complejo de Moore de su conjunto singular.
\end{remark}

\begin{example}[nervio de una categoría] Sea $\mathscr{C}$ una categoría localmente pequeña. Definimos el \textbf{nervio} de  $\mathscr{C}$ como el conjunto simplicial dado por los conjuntos
\begin{align*}
N(\mathscr{C})_n = \hom(\mathbf{n},\mathscr{C}) &= \{(f_1,\dots, f_k) : f_i \in \operatorname{mor} \mathscr{C}, \ \operatorname{cod} f_i = \operatorname{dom}f_{i+1}\}\\
& = \{x_0 \xrightarrow{f_1} x_1 \xrightarrow{f_2} \cdots \xrightarrow{f_n} x_n\}
\end{align*}
de $n$-uplas de morfismos componibles junto con los mapas 
\[
d_i(f_1, \dots, f_n) = (f_1, \dots,f_{i-1}, \ f_i \circ f_{i+1} \ ,\dots,f_n)
\]
y
\[
s_i(f_1, \dots, f_n) = (f_1, \dots, f_i, 1, f_{i+1},\dots,f_n),
\]
para cada $0 \leq i \leq n$. 

En particular, si $\mathscr{C} = BG$ es el grupoide asociado a un grupo $G$, entonces su nervio consiste de $n$-uplas de elementos de $G$ y los mapas de caras y degeneraciones están dados por
\[
d_i(g_1, \dots, g_n) = (g_1,\dots, g_{i-1},g_ig_{i+1},\dots,g_n), \quad s_i(g_1, \dots,g_n) = (g_1, \dots,g_i,1,g_{i+1},\dots,g_n).
\]
\end{example}

\begin{definition} Dados dos complejos simpliciales $X,Y : \ordcat^{op} \to \cat{Set}$, un \textbf{morfismo de conjuntos simpliciales} de $X$ a $Y$ es una transformación natural $f : X \to Y$. Concretamente, esto consiste en dar una familia de funciones $f_n : X_n \to Y_n$ tales que, para cada $0 \leq i \leq n$, los siguientes diagramas conmutan
\begin{center}
\begin{tikzpicture}
\matrix (m) [matrix of math nodes,row sep=3em,column sep=4em,minimum width=2em]
  {
     X_n & X_{n-1} & & X_n & X_{n+1}\\
     Y_n & Y_{n-1} & & Y_n & Y_{n+1}\\};
  \path[-stealth]
    (m-1-1) edge node [left] {$f_n$} (m-2-1)
    (m-1-4) edge node [left] {$f_n$} (m-2-4)
    (m-1-2) edge node [right] {$f_{n-1}$} (m-2-2)
    (m-1-5) edge node [right] {$f_{n+1}$} (m-2-5)
    (m-1-1) edge node [above] {$d_i$} (m-1-2)
    (m-2-1) edge node [below] {$d_i$} (m-2-2)
    (m-1-4) edge node [above] {$s_i$} (m-1-5)
    (m-2-4) edge node [below] {$s_i$} (m-2-5);
\end{tikzpicture} 
\end{center}

Es decir, un morfismo de conjuntos simpliciales consiste de una colección de aplicaciones que sea compatible con los mapas de caras y degeneraciones.
\end{definition}

\begin{remark} Los conjuntos simpliciales junto con los morfismos antes definidos forman una categoría\footnote{Esta es precisamente la categoría de prehaces de $\ordcat^{op}$.}, que notaremos $\cat{sSet}$.
\end{remark}

\begin{remark} Un morfismo $f : K \to L$ de complejos simlpiciales ordenados induce a su vez un morfismo de conjuntos simpliciales dado por $f_n[v_{i_1}, \dots, v_{i_n}] := [f(v_{i_1}), \dots, f(v_{i_n})]$ para cada $n$-símplex $[v_{i_1}, \dots, v_{i_n}]$ (posiblemente degenerado) de $K$.

De forma similar, una función continua $f : X \to Y$ induce un morfismo $f_* : \mathcal{S}(X) \to \mathcal{S}(Y)$ entre conjuntos singulares vía la postcomposición. De hecho, 
\end{remark}

\begin{definition} El \textbf{funtor singular} $\mathcal{S} : \cat{Top} \to \cat{sSet}$ asigna a cada espacio su conjunto singular, y a cada función continua $f : X \to Y$ el morfismo simplicial $f_* : \mathcal{S}(X) \to \mathcal{S}(Y)$ dado por $(f_*)_n(\sigma) =  f \circ \sigma$ para cada $\sigma: |\ss{n}| \to X$.
\end{definition}

\begin{remark} Si $X$ es un conjunto simplicial, el conjunto $X_n$ está determinado por los morfismos de conjuntos simpliciales de $\ss{n}$ a $X$. 

Concretamente, por el lema de Yoneda tenemos una biyección natural
\[
\hom_{\cat{sSet}}(\ss{n},X) \simeq X_n
\]

que a cada elemento $x \in X_n$ le asigna un morfismo de conjuntos simpliciales $\iota_x : \ss{n} \to X$ que satisface $\iota_x(1_{\ord{n}}) = x$. 

Esto se corresponde con la intuición que proveen los ejemplos anteriores, en los que los $n$-símplices de un conjunto simplicial son alguna «manifestación» de el $n$-símplex estándar: como la cara de un $m$-símplex de dimensión mayor, como un $m$-símplex de dimensión menor que represente un colapso del mismo, o como el $n$-símplex singular de un espacio topológico.\\
\end{remark}

\begin{definition} Sea $X$ un conjunto simplicial. Un $n$-simplex $x \in X_n$ se dice \textbf{degenerado} si existe $y \in X_{n-1}$ tal que $s_i(y) = x$ para algún $i \in \natzero{n}$. En caso contrario, decimos que $x$ es \textbf{no degenerado}.
\end{definition}

\section{Realización Geométrica}

Durante la materia vimos como a partir de un complejo simplicial $K$ podemos construir un espacio topológico $|K|$, la realización geométrica de $K$. Siguiendo esta idea, queremos extender esta noción al contexto de los conjuntos simpliciales.

\begin{definition}[realización geométrica, primera definición] Sea $X$ un conjunto simplicial. Dotando a cada conjunto $X_n$ de la topología discreta, definimos la \textbf{realización geométrica} de $X$ como el espacio
\[
|X| = \left(\coprod_{n \geq 0}X_n \times |\ss{n}|\right)\Big/\sim
\]
donde identificamos a los puntos $(x,|d^i|(p)) \sim (d_i(x),p)$ y $(x,|s^i|(p)) \sim (s_i(x),p)$.
\end{definition}

Esto formaliza la intuición anterior: para cada $x \in X_n$ construimos una copia del $n$-símplex estándar y los pegamos en función de si son caras o colapsos unos de otros.

Además, como veremos en breve esta asignación es funtorial, teniéndose así un análogo a la realización geométrica para complejos simpliciales.

\begin{proposition} La realización geométrica del $n$-símplex estándar $\ss{n} :\ordcat^{op} \to \cat{Set}$ es homeomorfa a la realización geométrica del $n$-símplex estándar como complejo simplicial.
\end{proposition}
\begin{proof} Para evitar ambiguedades, en toda la demostración notaremos $|\ss{n}|$ exclusivamente para referirnos a la realización geométrica del $n$-símplex como complejo simplicial. Por otro lado, notaremos $|\ordcat(-,n)|$.

Para cada $k \in \N_0$ tenemos una aplicación
\[
(f,x) \in \ordcat(k,n) \times |\Delta^k| \to |f|(x) \in |\Delta^n|,
\]
que resulta continua pues cada espacio $\ordcat(k,n)$ es discreto.

Ésta familia de funciones induce un morfismo en el coproducto que pasa al cociente por las identificaciones de la realización geométrica, pues si $g$ es un mapa de cocara o codegeneración, entonces
\[
r(g^*(f),x) = |fg|(x) = r(f,|g|(x)).
\]
Se tiene entonces una función continua $r : [(f,x)] \in |\ordcat(-,n)| \to |f|(x) \in |\Delta^n|$. 

Por otro lado, podemos considerar la inclusión $i : |\ss{n}| \to |\ordcat(-,n)|$ dada por la composición
\[
|\ss{n}| \xrightarrow{\simeq} \{id\} \times |\ss{n}| \hookrightarrow |\ordcat(-,n)|,
\]
que satisface $ri = 1$ pues
\[
ri(x) = [(id,x)] = |id|(x) = x
\]
para todo $x \in |\ss{n}|$. En particular sabemos que $i$ es inyectiva. Por lo tanto, para terminar basta ver que $i$ es sobreyectiva, en cuyo caso es un homeomorfismo con inversa $f$.

Equivalentemente, resta ver que todo elemento $[(f,x)]$ está relacionado con un punto de la forma $(id,y)$ para cierto $y \in |\ss{n}|$. En efecto, sea $(f,x) \in |\ordcat(-,n)|$ con $f : \ord{k} \to \ord{n}$ un morfismo de posets y $x \in |\ss{k}|$. Sabemos entonces que existen mapas de cocara o codegeneración $f_1, \dots, f_n$ tales que $f = f_1 \cdots f_n$. En consecuencia es
\begin{align*}
[(f,x)] &= [(f_1 \cdots f_n, x)] = [(f_n^* \circ \cdots \circ f_1^*(id), x)]\\ &= [(id,|f_1 \cdots f_n|(x))] = [(id, |f|(x))],
\end{align*}
lo que concluye la demostración.
\end{proof}

\begin{definition} Sea $X$ un conjunto simplicial. Su \textbf{categoría de símplices} es la categoría coma $\Delta \downarrow X$, donde $\Delta: \ordcat \to \cat{Set}^{\ordcat^{op}}$ es el embedding de Yoneda de $\ordcat$ y $X : \ordcat \to \cat{Set}^{\ordcat^{op}}$ es el funtor que vale constantemente $X$. Concretamente, los objetos de $\Delta \downarrow X$ son \textit{símplices} de $X$, entendidos como morfismos simpliciales $\ss{n} \to X$, y las flechas son morfismos $\theta : \ord{n} \to \ord{m}$ de posets tales que el diagrama
\begin{center}
\begin{tikzpicture}
\matrix (m) [matrix of math nodes,row sep=3em,column sep=4em,minimum width=2em]
  {
     \ss{n}& &\\
     & & X\\
     \ss{m} & &\\};
  \path[-stealth]
    (m-1-1) edge node [above] {$\sigma$} (m-2-3)
    (m-1-1) edge node [left] {$\theta_*$} (m-3-1)
    (m-3-1) edge node [below] {$\tau$} (m-2-3);
\end{tikzpicture} 
\end{center}
conmuta, donde $\theta_*$ es la postcomposición por $\theta$.
\end{definition}

\begin{theorem} Si $X$ es un conjunto simplicial, entonces
\[
|X| \simeq \catcolim{\substack{\ss{n} \to X \\ \text{ en $\Delta \downarrow X$}}}{|\ss{n}|}.
\]
\end{theorem}
\begin{proof} Observemos que si $Z$ es un espacio topológico arbitrario, una función continua $g : |X| \to Z$ se corresponde unívocamente con una función continua $\tilde{g} : \coprod_{n \geq 0} X_n \times |\ss{n}| \to Z$ que sea compatible con la relación que identifica caras y colapsos.

A su vez, usando la propiedad universal del coproducto y el hecho de que cada espacio $X_n$ tiene la topología discreta, esto equivale a dar funciones
\[
\lambda_x : |\ss{n}| \to Z
\]
para cada $x \in X_n$ y $n \geq 0$, que satisfagan las condiciones de compatibilidad de antes: esto es, que para cada morfismo de posets\footnote{Si bien la condición original era sobre los mapas de caras y degeneración, que son imagen por $X$ de los mapas de cocaras y codegeneración en $\ordcat$. Al clausurar la relación de equivalencia se tiene una condición equivalente reemplazando éstos por cualquier función $X(\theta)$ con $\theta$ un morfismo en $\ordcat$.} $\theta : \ord{n} \to \ord{m}$ se tenga $\lambda_{X(\theta)(x)} = \lambda_x \circ |\theta| = |\theta^*|(\lambda_x)$.

Recordemos ahora que, por el lema de Yoneda, tenemos que
\[
\hom(\ss{n},\ss{m}) \simeq \ordcat(n,m) = \{\theta : \ord{n} \to \ord{m} : \text{ $\theta$ es morfismo de posets}\}
\]
y
\[
\hom(\ss{n}, X) \simeq X_n
\]
para cada $n,m \geq 0$. Por lo tanto, cada elemento $x \in X_n$ se corresponde a un único morfismo simplicial $\sigma_x : \ss{n} \to X$, y toda flecha $\ss{n} \to \ss{m}$ es la poscomposición por cierto morfismo de posets $\theta: \ord{n} \to \ord{m}$.

En éstos terminos, la información anterior se puede describir como un morfismo $\lambda_\sigma : |\ss{n}| \to Z$ para cada símplex $\sigma : \ss{n} \to X$, de forma que para todo morfismo de posets $\theta : \ord{n} \to \ord{m}$ y símplices $\sigma : \ss{n} \to X, \tau : \ss{m} \to X$ tales que $\sigma = \tau \  \theta_*$ se tenga que $\lambda_\sigma = \lambda_\tau |\theta|$. 

Es decir, dar un morfismo $g : |X| \to Z$ es equivalente a dar un cocono sobre $Z$ para el funtor $F : \Delta \downarrow X \to \mathsf{Top}$ que envía $(\sigma : \ss{n} \to X) \xrightarrow{\theta_*} (\tau : \ss{m} \to X)$ a $|\ss{n}| \xrightarrow{|\theta|} |\ss{m}|$.

Por otro lado, tenemos un cocono sobre $X$ dado por las funciones
\[
\iota_\sigma : p \in |\ss{n}| \to [(x,p)] \in |X|,
\]
para cada $\sigma : \ss{n} \to X$ que está en correspondencia con $x \in X_n$. Por las observaciones anteriores, sabemos que $(\lambda_\sigma)_{\sigma \in \Delta \downarrow X}$ se factoriza por $(\iota_\sigma)_{\sigma \in \Delta \downarrow X}$ a través de $g$ ya que es
\[
\lambda_\sigma(p) = g([(x,p)]) = g(\iota_\sigma(x,p))
\]
para todo $\sigma : \ss{n} \to X$ y $p \in |\ss{n}|$. Pero como notamos anteriormente, de existir una función continua $|X| \to Z$ está determinada por lo que vale en la imagen de cada morfismo $\iota_\sigma$. 
Hemos visto entonces que todo cono sobre $F : \Delta \downarrow X \to \mathsf{Top}$ se factoriza a través de $(\iota_\sigma)_\sigma$ de forma única, y esto es precisamente que
\[
|X| \simeq \catcolim{\substack{\ss{n} \to X \\ \text{ en $\Delta \downarrow X$}}}{|\ss{n}|}.
\]
\end{proof}

\begin{proposition} Se tiene un funtor
\[
| \cdot | : \cat{sSet} \to \cat{Top}
\]
que asigna a cada conjunto simplicial $X$ su realización geométrica, y a cada morfismo simplicial $f : X \to Y$ una flecha $|f| : |X| \to |Y|$ inducida por cada función $f_n \times 1 : X_n \times |\ss{n}| \to Y_n \times |\ss{n}|$.
\end{proposition}
\begin{proof} En primer lugar, notemos que las aplicaciones $\{f_n \times 1\}_{n \geq 1}$ inducen una función continua
\begin{center}
\begin{tikzcd}
 \coprod_{n \geq 0} X_n \times |\ss{n}| \arrow{r}{\widetilde{f}} &  \coprod_{n \geq 0} Y_n \times |\ss{n}|\\
 X_n \times |\ss{n}| \arrow[hookrightarrow]{u} \arrow{r}{f_n \times 1} & Y_n \times |\ss{n}| \arrow[hookrightarrow]{u}
\end{tikzcd}
\end{center}

entre los coproductos. Como $f$ es un morfismo simplicial, si $\theta : \ord{m} \to \ord{m}$ es un morfismo de posets entonces el diagrama
\begin{center}
\begin{tikzcd}
 X_n \times |\ss{n}| \arrow{r}{X\theta}\arrow{d}[left]{f_n \times 1} &  X_m \times |\ss{n}| \arrow{d}{f_m \times 1}\\
 Y_n \times |\ss{n}| \arrow{r}{Y \theta} & Y_m \times |\ss{n}| 
\end{tikzcd}
\end{center}
conmuta. Por lo tanto, se tiene que
\[
\widetilde{f}(X\theta(x),p) = (f_mX\theta(x),p) = (Y\theta f_n(x),p)
\]
y
\[
\widetilde{f}(x,|\theta|(p)) = (f_n(x),|\theta|(p)),
\]
lo que nos dice que $\widetilde{f}$ manda puntos relacionados en puntos relacionados. En consecuencia, está bien definida la función continua $|f| : |X| \to |Y|$ que envía $[(x,p)]$ a $[(f_n(x),p)]$ si $x \in X_n$, y de esta caracterización se deduce la funtorialidad de $|\cdot|$.
\end{proof}

\section{La adjunción entre la realización geométrica y el funtor singular}

Como vimos en la sección anterior, la realización geométrica está directamente relacionada con el conjunto singular de un espacio topológico. El teorema que sigue describe explícitamente esta relación en terminos de los functores $| \cdot |$ y $\mathcal{S}$.

\begin{theorem} Existe una adjunción
\begin{center}
\begin{tikzcd}
    \cat{Top} \arrow[bend left=35]{rr}{| \  \cdot \ |}
  & \rotatebox{90}{$\vdash$}
  & \cat{sSet} \arrow[bend left=35]{ll}{\mathcal{S}}
\end{tikzcd}
\end{center}
En otras palabras, se tiene una biyección
\[
\cat{Top}(|X|,Y) \simeq \cat{sSet}(X,\mathcal{S}(Y))
\]
entre las funciones continuas $|X| \to Y$ y los morfismos simpliciales $X \to \mathcal{S}(Y)$. Más aún, esta biyección es natural tanto en $X$ como en $Y$.
\end{theorem}
\begin{proof} Recordemos que hay una correspondencia biyectiva entre morfimos simpliciales $\sigma: \ss{n} \to X$ y símplices $x \in X_n$. Notamos $\sigma_x$ al morfismo simplicial asociado a un símplex $x$. Vimos también que una función continua $g : |X| \to Y$ se corresponde unívocamente con una colección de funciones continuas $g_x : \ss{n_x} \to Y$ tales que 
\[
g([(x,p)]) = g_x(p)
\]
para todo símplex $x \in X_n$ y $p \in |\ss{n}|$. 

En vista de esto, definimos la aplicación
\begin{align*}
e : \cat{Top}(|X|,&Y) \to \cat{sSet}(X,\mathcal{S}(Y))\\
& g \longmapsto e(g)(x) := g_x
\end{align*}
Como $g$ es un morfismo cuyo dominio es $|X|$, las funciones $g_x$ deben satisfacer
\[
g_{X\theta(x)} = |\theta|^*(g_x)
\]
para todo símplice $x$ y morfismo de posets $\theta$. En particular, esto prueba que $e(g)$ es efectivamente un morfismo de complejos simpliciales, por lo que $e$ está bien definida.

Por otro lado, definimos otra aplicación
\begin{align*}
\rho : \cat{sSet}(X,\mathcal{S}(&Y)) \to \cat{Top}(|X|,Y)\\
& f \longmapsto \rho(f)
\end{align*}
del siguiente modo: dado un morfismo simplicial $f$, para cada $f_n : X_n \to \mathcal{S}(Y)_n$ definimos las funciones
\[
\widetilde{f_n} : (x,p) \in X_n \times |\ss{n}| \mapsto f_n(x)(p) \in Y,
\]
las cuales inducen a su vez una función continua
\[
\coprod_{n \geq 0} \widetilde{f_n} : \coprod_{n \geq 0} X_n \times |\ss{n}| \to Y.
\]

Como $f$ es un morfismo simplicial, para cada $\theta \in \ordcat(m,n)$ se tiene que
\[
\widetilde{f_m}(X\theta(x), p) = (f_mX\theta)(x)(p) = |\theta|^*(f_n(x))(p) = (f_n(x) |\theta|)(p) = \widetilde{f_n}(x,|\theta|(p)).
\]
Esto dice que la función $\coprod_{n \geq 0} \widetilde{f_n}$ pasa al cociente por las identificaciones de la realización geométrica, y por lo tanto, induce una aplicación continua $|X| \to Y$ que tomamos como imagen de $f$ por $\rho$. Concretamente, dados $x \in X_n$ y $p \in |\ss{n}|$ es
\[
\rho(f)([(x,p)]) = f_n(x)(p).
\]

Veamos ahoar que $e$ y $\rho$ son aplicaciones inversas. Si tomamos $f : X \to \mathcal{S}(Y)$ simplicial, es
\[
e\rho(f)(x)(p) = \rho(f)([(x,p)]) = f_n(x)(p)
\]
para todo $x \in X_n$ y $p \in |\ss{n}|$. Por lo tanto, se tiene que $e \rho(f) = f$ y, en general, $e \rho = 1$.

Recíprocamente, si $g : |X| \to Y$ es una función continua, entonces para cada $(x,p) \in X_n \times |\ss{n}|$ es
\[
\rho e(g)([(x,p)]) = e(g)_n(x)(p) = g_x(p) = g([(x,p)]),
\]
lo cual prueba que $\rho e  = 1$.

Para terminar, veamos que el isomorfismo es natural tanto en $X$ como $Y$. Fijemos $g : |X| \to Y$ continua y $(x,p) \in X_n \times |\ss{n}|$. Dado un morfismo simplicial $f : X' \to X$, debemos ver que el diagrama
\begin{center}
\begin{tikzcd}
\cat{Top}(|X|,Y) \arrow{r}{e}\arrow{d}[left]{|f|^*} \  &\ \cat{sSet}(X,\mathcal{S}(Y))\arrow{d}{f^*}\\ 
\cat{Top}(|X'|,Y)\arrow{r}{e} & \cat{sSet}(X',\mathcal{S}(Y)) 
\end{tikzcd}
\end{center}
conmuta. Por un cálculo directo es
\begin{align*}
f^*(e(g))(x)(p) &= e(g)(f(x))(p) = g([(f(x),p)])\\
& = g|f|([(x,p)]) = e(g|f|)(x)(p)\\
& = e(|f|^*(g))(x)(p),
\end{align*}
y por lo tanto $f^*e = e|f|^*$. Del mismo modo, si $f : Y \to Y'$ es continua entonces
\begin{center}
\begin{tikzcd}
\cat{Top}(|X|,Y) \arrow{r}{e}\arrow{d}[left]{f_*} \  &\ \cat{sSet}(X,\mathcal{S}(Y))\arrow{d}{f_*}\\ 
\cat{Top}(|X|,Y')\arrow{r}{e} & \cat{sSet}(X,\mathcal{S}(Y')) 
\end{tikzcd}
\end{center}
conmuta pues
\begin{align*}
f_*(e(g))(x)(p) &= f(e(g)(x)(p)) = fg([(x,p)])\\
&= e(fg)(x)(p) = e(f_*(g))(x)(p).
\end{align*}
\end{proof}

\begin{remark} En \cite{GJ} se encuentra una demostración «puramente categórica» del teorema anterior usando que todo prehaz es un colímite de funtores representables, y que las realizaciones geométricas son colímites de símplices (topológicos) estándar.

Concretamente, usando la relacion entre funtores $\hom$ y colímites (ver \cite{context}, sección 3.4) es
\begin{align*}
\cat{Top}(|X|,Y) &\simeq \catcolim{\substack{\ss{n} \to X \\ \text{ en $\Delta \downarrow X$}}}{\cat{Top}(|\ss{n}|,Y)}\\ &\simeq \catcolim{\substack{\ss{n} \to X \\ \text{ en $\Delta \downarrow X$}}}{\cat{sSet}(\ss{n},\mathcal{S}(Y))}\\&\simeq  \cat{sSet}(X,\mathcal{S}(Y)).
\end{align*}
\end{remark}

\begin{thebibliography}{}
\bibitem{GJ} P. Goerss y J. Jardine, \textit{Simplicial Homotopy Theory}. Birkhäuser, 2010.
\bibitem{F} G. Friedman. \textit{An elementary introduction to simplicial sets}, arXiv:0809.4221v5 [at], 2016.
\bibitem{context} E. Riehl. \textit{Category Theory in Context}. Dover, 2016.
\end{thebibliography}
\end{document}


